\documentclass{article}
\usepackage{amsmath, amssymb, amsthm}

\title{IMC 2024 \\ Second Day, August 8, 2024}
\date{}

\begin{document}
\maketitle

\section*{Problem 6}
Prove that for any function $f: \mathbb{Q} \to \mathbb{Z}$, there exist $a, b, c \in \mathbb{Q}$ such that $a < b < c$, $f(b) \geq f(a)$, and $f(b) \geq f(c)$.

(10 points)

\section*{Problem 7}
Let $n$ be a positive integer. Suppose that $A$ and $B$ are invertible $n \times n$ matrices with complex entries such that $A + B = I$ and
\[
(A^2 + B^2)(A^4 + B^4) = A^5 + B^5.
\]
Find all possible values of $\det(AB)$ for the given $n$.

(10 points)

\section*{Problem 8}
Define the sequence $x_1, x_2, \ldots$ by $x_1 = 2$, $x_2 = 4$, and
\[
x_{n+2} = 3 x_{n+1} - 2 x_n + \frac{2^n}{x_n} \quad \text{for } n \geq 1.
\]
Prove that $\lim_{n \to \infty} \frac{x_n}{2^n}$ exists and satisfies
\[
\frac{1 + \sqrt{3}}{2} \leq \lim_{n \to \infty} \frac{x_n}{2^n} \leq \frac{3}{2}.
\]
(10 points)

\section*{Problem 9}
A matrix $A = (a_{ij})$ is called \emph{nice} if:
\begin{itemize}
\item[(i)] the set of all entries of $A$ is $\{1, 2, \ldots, 2t\}$ for some integer $t$;
\item[(ii)] $a_{i,j} \leq a_{i,j+1}$ and $a_{i,j} \leq a_{i+1,j}$;
\item[(iii)] if $a_{i,j} = a_{k,\ell}$, then either $i = k$ or $j = \ell$;
\item[(iv)] for each $s = 1, \ldots, 2t-1$, there exist $i \neq k$ and $j \neq \ell$ with $a_{i,j} = s$ and $a_{k,\ell} = s+1$.
\end{itemize}
Prove that the number of nice $m \times n$ matrices is even. For example, the only two nice $2 \times 3$ matrices are $\begin{pmatrix} 1 & 1 & 1 \\ 2 & 2 & 2 \end{pmatrix}$ and $\begin{pmatrix} 1 & 1 & 3 \\ 2 & 4 & 4 \end{pmatrix}$.

(10 points)

\section*{Problem 10}
We say that a square-free positive integer $n$ is \emph{almost prime} if
\[
n \mid x^{d_1} + x^{d_2} + \ldots + x^{d_k} - kx
\]
for all integers $x$, where $1 = d_1 < d_2 < \ldots < d_k = n$ are all positive divisors of $n$. Suppose $r = 2^{2^m}+1$ is a Fermat prime, $p$ is a prime divisor of an almost prime $n$, and $p \equiv 1 \pmod{r}$. Show $d_i \equiv 1 \pmod{r}$ for all $i$.

(10 points)

\end{document}
